% Werte und Aufgabenstellung: aufgaben/Blatt07.pdf, S. 2.
\subsection*{Gegeben und Modell}
Der schwimmende Eisberg habe vor dem Schmelzen das Gesamtvolumen $V$ und die
Dichte $\rho_\mathrm E=\qty{920}{\kilogram\per\metre\cubed}$. Für das Meer
ist $\rho_\mathrm M=\qty{1030}{\kilogram\per\metre\cubed}$ gegeben. Die
Luftdichte und damit der Luftauftrieb werden gemäss Blatt vernachlässigt.
Die Meeresoberfläche $A_\mathrm M$ sei konstant.

\subsection*{a) Volumen oberhalb der Oberfläche}
Sei $V_\mathrm{u}$ das eingetauchte Volumen. Im Gleichgewicht ist die
Gewichtskraft des Eises gleich dem Auftrieb des verdrängten Meerwassers:
\begin{equation}
  \rho_\mathrm E Vg=\rho_\mathrm M V_\mathrm{u}g,
  \qquad \frac{V_\mathrm{u}}{V}=\frac{\rho_\mathrm E}{\rho_\mathrm M}.
\end{equation}
Das Volumen oberhalb der Oberfläche ist $V-V_\mathrm{u}$. Daher
\begin{ergebnis}[a) Oberhalb der Oberfläche]
\begin{equation*}
  \frac{V-V_\mathrm{u}}{V}
  =1-\frac{920}{1030}\approx \num{0.1068}=\qty{10.68}{\percent}.
\end{equation*}
\end{ergebnis}
Die übrigen rund \qty{89.32}{\percent} sind eingetaucht. Diese Anteile
beziehen sich auf Volumina, nicht auf Höhen oder Längen.

\subsection*{b) Änderung des Meeresspiegels}
Die Masse $m=\rho_\mathrm E V$ bleibt beim Schmelzen erhalten. Die Dichte
$\rho_\mathrm S$ des Schmelzwassers ist auf dem Blatt nicht angegeben. Wir
nehmen an, dass der Eisberg aus praktisch salzfreiem Eis besteht und zu
Süsswasser schmilzt. Dann ist das Schmelzwasservolumen $m/\rho_\mathrm S$,
nicht automatisch $m/\rho_\mathrm M$.

Wir modellieren das Meerwasser als inkompressibel mit unveränderter Dichte
$\rho_\mathrm M$ und nehmen additive Volumina ohne Mischungs- oder
Temperaturkontraktion an. Das Schmelzwasser kann sich verteilen, ohne dass
sein beitragendes Volumen durch Mischung verändert wird. Die reale
Salzgehaltsänderung und ihre Wirkung auf die Dichte sind damit nicht erfasst.
Vor dem Schmelzen belegt der Eisberg unter der Oberfläche das Volumen
$V_\mathrm{u}=m/\rho_\mathrm M$. Anschliessend entfällt diese Verdrängung,
und das Schmelzwasser nimmt das Volumen $m/\rho_\mathrm S$ ein. Bei
konstanter Oberfläche folgt aus $A_\mathrm M h=\Delta V$:
\begin{ergebnis}[b) Meeresspiegeländerung]
\begin{equation*}
  h=\frac{\rho_\mathrm E V}{A_\mathrm M}
    \left(\frac{1}{\rho_\mathrm S}-\frac{1}{\rho_\mathrm M}\right).
\end{equation*}
\end{ergebnis}
Für $\rho_\mathrm S<\rho_\mathrm M$ ist $h>0$. Ein Zahlenkoeffizient ist
nicht allein aus den Angaben des Blattes bestimmbar. Als ausdrücklich
zusätzliche Näherung wählen wir $\rho_\mathrm S=\qty{1000}{\kilogram\per\metre\cubed}$;
damit ergibt sich
\begin{equation}
  h=\left(\frac{920}{1000}-\frac{920}{1030}\right)\frac{V}{A_\mathrm M}
    \approx \num{0.02680}\,\frac{V}{A_\mathrm M}.
\end{equation}
Ohne konkrete Angaben zu $V$ und $A_\mathrm M$ folgt keine absolute Höhe.
Falls die Flüssigkeiten vollständig durchmischt werden und dabei ihr Volumen
ändern, benötigt man stattdessen die Enddichte beziehungsweise ein
Mischungsgesetz; das Ergebnis oben ist dann nicht ohne Weiteres exakt.

\begin{kontrolle}[Einheit und Grenzfall]
$V/A_\mathrm M$ hat die Einheit einer Länge. Für gleiche Dichten von
Schmelzwasser und Meerwasser verschwindet $h$: Das Schmelzwasser ersetzt dann
genau das zuvor verdrängte Volumen. Dieser Sonderfall darf für Süsswasser in
Meerwasser nicht ungeprüft vorausgesetzt werden.
\end{kontrolle}

\begin{merksatz}
Bei schwimmenden Körpern bestimmt die Masse die Verdrängung. Für die
Pegeländerung nach einem Phasenwechsel ist zusätzlich das Endvolumen des
Materials und gegebenenfalls das Mischungsverhalten entscheidend.
\end{merksatz}
